\documentclass{beamer}

% Theme selection
\usetheme{Madrid}
\usecolortheme{whale}

% Packages
\usepackage[utf8]{inputenc}
\usepackage[T1]{fontenc}
\usepackage{graphicx}
\usepackage{booktabs}
\usepackage{listings}
\usepackage{xcolor}

% Technical listings style
\lstset{
    basicstyle=\ttfamily\footnotesize,
    breaklines=true,
    captionpos=b,
    commentstyle=\color{gray},
    keywordstyle=\color{blue},
    stringstyle=\color{red}
}

% Title information
\title[TabuLLM Architecture]{TabuLLM: Synthetic Tabular Data Generation}
\subtitle{Architecture, Challenges, and Future Directions}
\author{Development Team}
\date{\today}

\begin{document}

% --- SLIDE 1: Title ---
\begin{frame}
    \titlepage
\end{frame}

% --- SLIDE 2: Introduction ---
\begin{frame}{What is TabuLLM?}
    \begin{itemize}
        \item \textbf{Core Objective}: A tool to generate synthetic tabular data and augment existing datasets using Large Language Models (LLMs).
        \item \textbf{Key Capabilities}:
        \begin{itemize}
            \item Interactive UI for dataset manipulation.
            \item Automated metadata extraction and EDA.
            \item Smart context-aware generation (filling specific rows or generating new ones).
        \end{itemize}
        \item \textbf{Stack}: Python-based, utilizing Dash for the frontend and variety of LLM providers (Local \& Cloud).
    \end{itemize}
\end{frame}

% --- SLIDE 3: High-Level Architecture ---
\begin{frame}{High-Level Architecture}
    \begin{description}
        \item[Frontend (UI)] Built with \textbf{Plotly Dash}, \textbf{Bootstrap}, and \textbf{AG-Grid} for high-performance table interactions.
        \item[Core Engine] Manages application state, background callbacks (via Diskcache), and logging.
        \item[Data Services] Handles file I/O (CSV/Excel), EDA metrics, and context selection algorithms (Random, Dense, Diverse).
        \item[LLM Integration] Multi-provider client interface supporting Local OpenAI-compatible servers (LM Studio) and Google Gemini.
    \end{description}
\end{frame}

% --- SLIDE 4: LLM Service: Optimization & Robustness ---
\begin{frame}{LLM Service: Optimization \& Robustness}
    \begin{itemize}
        \item \textbf{Token Optimization}:
        \begin{itemize}
            \item Uses a \texttt{Compact JSON} format (list-of-lists) to minimize overhead compared to standard JSON objects.
            \item Strategic chunking of row payloads to stay within context windows.
        \end{itemize}
        \item \textbf{Syntactic Reliability}:
        \begin{itemize}
            \item \textbf{Robust JSON Parsing}: Multi-step cleaning of raw LLM output.
            \item \textbf{Self-Repair Loop}: If JSON is malformed, a secondary "repair prompt" focuses specifically on fixing the syntax error coordinates.
        \end{itemize}
        \item \textbf{Resilience}: Exponential backoff for rate-limited APIs (429 errors).
    \end{itemize}
\end{frame}

% --- SLIDE 5: Data Management: Smart Context ---
\begin{frame}{Data Management: Smart Context Selection}
    Giving the LLM the "right" examples is more critical than giving it "many" examples.
    \vspace{0.2cm}
    \begin{block}{Selection Strategies}
        \begin{itemize}
            \item \textbf{Random}: Baseline sampling for general distribution.
            \item \textbf{Dense}: Prioritizes rows with the most non-null values to maximize information density.
            \item \textbf{Diverse (K-Means)}: Uses clustering on normalized numerical features to select archetypal rows that cover the full feature space.
        \end{itemize}
    \end{block}
\end{frame}

% --- SLIDE 6: Difficulties Faced ---
\begin{frame}{Difficulties Faced during Development}
    \begin{itemize}
        \item \textbf{LLM Hallucinations}: Ensuring numeric types (ints vs. floats) and specific formatting (no currency symbols in raw data).
        \item \textbf{Latency vs. Quality}: Balancing the speed of local 4B/7B models against the superior quality and instructions-following of large cloud models (Gemini 1.5).
        \item \textbf{Context Management}: Tabular data is verbose. Estimating token counts dynamically helps users manage cost and avoid truncation.
        \item \textbf{JSON Fragility}: Even sophisticated models often fail to close brackets or comma separators when generating hundreds of rows.
    \end{itemize}
\end{frame}

% --- SLIDE 7: Future Challenges ---
\begin{frame}{Future Challenges \& Roadmap}
    \begin{enumerate}
        \item \textbf{Scalability}: Transitioning from in-memory Pandas/AG-Grid to streaming architectures for datasets with millions of rows.
        \item \textbf{Semantic Constraints}: Implementing "Cross-Column Guards" (e.g., if Age < 18, then Occupation $\neq$ "Doctor").
        \item \textbf{Privacy}: Integrating Differential Privacy (DP) mechanisms to ensure LLM-generated data doesn't leak sensitive source records.
        \item \textbf{Relational Structure}: Handling multi-table datasets with foreign keys via graph-based context.
    \end{enumerate}
\end{frame}

% --- SLIDE 8: Conclusion ---
\begin{frame}{Conclusion}
    \begin{itemize}
        \item TabuLLM bridges the gap between raw LLM power and practical tabular data workflows.
        \item \textbf{Current State}: Functional prototyping tool with robust error recovery and smart sampling.
        \item \textbf{Next Steps}: Deepening semantic validation and improving scalability for enterprise-grade datasets.
    \end{itemize}
    \vspace{0.8cm}
    \centering
    \textbf{Questions?}
\end{frame}

\end{document}
